\documentclass[11pt,a4paper]{article}

\usepackage{fontawesome5}
\providecommand{\faPencilAlt}{\faIcon{pencil-alt}}

\usepackage{styles/couleurs}
\usepackage{styles/mise_en_page}
\usepackage{styles/boites}
\usepackage{styles/maths}
\usepackage{styles/figures}

\usepackage{enumitem}
\usepackage{array}
\usepackage{tabularx}
\usepackage{amsmath}
\usepackage{amssymb}
\usepackage{tikz}
\usepackage{xparse}
\usepackage{graphicx}
\usepackage{ragged2e}
\usepackage{microtype}

\usetikzlibrary{calc,positioning}

\setlength{\parindent}{0pt}
\setlength{\parskip}{0pt}
\raggedbottom

\setlist[itemize]{
  leftmargin=1.25em,
  itemsep=0.05em,
  topsep=0.1em
}

% =========================================================
% IMAGES
% Toutes les images sont dans le dossier images/
% =========================================================

\newcommand{\copieimage}[2][]{%
  \IfFileExists{images/#2.png}{%
    \includegraphics[#1]{images/#2.png}%
  }{%
    \IfFileExists{images/#2.jpg}{%
      \includegraphics[#1]{images/#2.jpg}%
    }{%
      \IfFileExists{images/#2.jpeg}{%
        \includegraphics[#1]{images/#2.jpeg}%
      }{%
        \fbox{%
          \parbox[c][3cm][c]{0.9\linewidth}{%
            \centering\sffamily\small
            Image introuvable :\\
            \texttt{\detokenize{#2}}
          }%
        }%
      }%
    }%
  }%
}

% Cadre léger autour des scans.
% Aucun recadrage : on conserve l'image entière.
\newcommand{\copieencadree}[2][]{%
  \begingroup
  \setlength{\fboxsep}{1.5pt}%
  \setlength{\fboxrule}{0.35pt}%
  \fcolorbox{black!28}{white}{%
    \copieimage[#1]{#2}%
  }%
  \endgroup
}

% =========================================================
% TITRE
% =========================================================

\newcommand{\titreAP}[1]{%
\begin{center}
  {\sffamily\LARGE\bfseries\color{TitreC}
  AP -- J'écris des maths}\\[-0.08em]

  {\sffamily\large\bfseries\color{SavoirC}
  #1}
\end{center}

\vspace{-0.35em}
}

% =========================================================
% BOUSSOLE
% =========================================================

\newcommand{\boussole}{%
\begin{tcolorbox}[
  enhanced,
  colback=SavoirBg,
  colframe=SavoirC!42,
  boxrule=0.45pt,
  arc=1.5mm,
  left=2mm,
  right=2mm,
  top=0.8mm,
  bottom=0.8mm,
  before skip=0.4mm,
  after skip=1.3mm
]

\small

{\sffamily\bfseries\color{SavoirC}
Si cela t'aide :}
\qquad
\textbf{\(+\)} quelque chose qui m'aide
\qquad
\textbf{\(?\)} quelque chose qui m'interroge
\qquad
\textbf{\(\rightarrow\)} quelque chose que je modifierais

\end{tcolorbox}
}

% =========================================================
% CORRECTION ÉCLAIR
% =========================================================

\newcommand{\correctioneclair}[1]{%
\begin{tcolorbox}[
  enhanced,
  colback=NeutreBg,
  colframe=NeutreC!35,
  boxrule=0.45pt,
  arc=1.5mm,
  borderline west={1mm}{0pt}{NeutreC},
  left=2.3mm,
  right=2.3mm,
  top=1mm,
  bottom=1mm,
  before skip=0.8mm,
  after skip=1.2mm
]

{\small
\textbf{Correction éclair :}\quad #1
}

\end{tcolorbox}
}

% =========================================================
% MISSION OUVERTE
% =========================================================

\newcommand{\missionouverte}[1]{%
\begin{tcolorbox}[
  enhanced,
  colback=ActionC!5,
  colframe=ActionC!48,
  boxrule=0.5pt,
  arc=1.5mm,
  left=2.4mm,
  right=2.4mm,
  top=1mm,
  bottom=1mm,
  before skip=0.8mm,
  after skip=1.4mm
]

{\sffamily\bfseries\small\color{ActionC}
Mission}
\quad
{\small #1}

\end{tcolorbox}
}

% =========================================================
% ZONE DE REMARQUES
% =========================================================

\newcommand{\remarques}[1]{%
\begin{tcolorbox}[
  enhanced,
  colback=white,
  colframe=NeutreC!28,
  boxrule=0.45pt,
  arc=1.5mm,
  left=2.5mm,
  right=2.5mm,
  top=1.5mm,
  bottom=1.5mm,
  height=#1,
  before skip=1.2mm,
  after skip=0.6mm
]

{\sffamily\bfseries\small\color{NeutreC}
Mes remarques}

\end{tcolorbox}
}

% =========================================================
% ÉTIQUETTES
% =========================================================

\newcommand{\etiquette}[1]{%
  {\sffamily\bfseries\small\color{TitreC}#1}
}

% =========================================================
% DOCUMENT
% =========================================================

\begin{document}

% =========================================================
% PAGE 1
% GÉOMÉTRIE : PRODUCTIONS 1 ET 3
% =========================================================

\titreAP{Problème de géométrie -- productions 1 et 3}

\boussole

\noindent
\begin{minipage}[c]{0.60\linewidth}

{\sffamily\bfseries\small\color{TitreC}
Énoncé}

\smallskip

\small

\(ABC\) est un triangle rectangle en \(B\) tel que
\(\widehat{BAC}=50^\circ\) et \(AC=4\) cm.

\(BCD\) est un triangle équilatéral.

\(CDE\) est un triangle tel que
\(\widehat{DCE}=80^\circ\) et \(CE=2{,}4\) cm.

\textbf{Calculer la longueur \(AE\).}

\end{minipage}
\hfill
\begin{minipage}[c]{0.37\linewidth}

\centering

\copieencadree[
  width=0.96\linewidth,
  height=0.15\textheight,
  keepaspectratio
]{problème de longueur figure}

\end{minipage}

\correctioneclair{%
\(\widehat{ACB}=40^\circ\) et
\(\widehat{BCD}=60^\circ\).
Donc
\(40^\circ+60^\circ+80^\circ=180^\circ\) :
\(A\), \(C\) et \(E\) sont alignés.
Ainsi
\(AE=4+2{,}4=6{,}4\) cm.
}

\missionouverte{%
Observe les deux productions.
Annote-les librement et note ce que tu remarques.
}

% ---------- Production 1 ----------

\etiquette{Production 1}

\vspace{0.2em}

\begin{center}

\copieencadree[
  width=0.94\linewidth,
  height=0.12\textheight,
  keepaspectratio
]{production 1}

\end{center}

\vspace{0.4em}

% ---------- Production 3 ----------

\etiquette{Production 3}

\vspace{0.2em}

\noindent
\begin{minipage}[c]{0.49\linewidth}
\centering

\copieencadree[
  width=0.95\linewidth,
  height=0.22\textheight,
  keepaspectratio
]{production 3 (1)}

\end{minipage}
\hfill
\begin{minipage}[c]{0.49\linewidth}
\centering

\copieencadree[
  width=0.95\linewidth,
  height=0.22\textheight,
  keepaspectratio
]{production 3 (2)}

\end{minipage}

\remarques{2.7cm}


% =========================================================
% PAGE 2
% GÉOMÉTRIE : PRODUCTION 4
% =========================================================

\newpage

\titreAP{Problème de géométrie -- production 4}

\boussole

\correctioneclair{%
\(\widehat{ACB}=40^\circ\),
\(\widehat{BCD}=60^\circ\) et
\(40^\circ+60^\circ+80^\circ=180^\circ\).
Donc \(A,C,E\) sont alignés et
\(AE=6{,}4\) cm.
}

\missionouverte{%
Observe cette production.
Annote-la et écris ce que tu en penses.
}

\begin{center}

\etiquette{Production 4 -- brouillon}

\vspace{0.3em}

\copieencadree[
  width=0.92\linewidth,
  height=0.58\textheight,
  keepaspectratio
]{production 4}

\end{center}

\remarques{3.3cm}


% =========================================================
% PAGE 3
% GÉOMÉTRIE : PRODUCTION 5
% =========================================================

\newpage

\titreAP{Problème de géométrie -- production 5}

\boussole

\correctioneclair{%
\(\widehat{ACB}=40^\circ\),
\(\widehat{BCD}=60^\circ\),
puis
\(40^\circ+60^\circ+80^\circ=180^\circ\).
Donc \(A,C,E\) sont alignés et
\(AE=6{,}4\) cm.
}

\missionouverte{%
Observe et annote cette production.
Écris ensuite tes remarques.
}

\begin{center}

\etiquette{Production 5}

\vspace{0.3em}

\copieencadree[
  width=0.93\linewidth,
  height=0.31\textheight,
  keepaspectratio
]{production 5 (1)}

\vspace{0.5em}

\copieencadree[
  width=0.93\linewidth,
  height=0.31\textheight,
  keepaspectratio
]{production 5 (2)}

\end{center}

\remarques{2.7cm}


% =========================================================
% PAGE 4
% GÉOMÉTRIE : PRODUCTION 6
% =========================================================

\newpage

\titreAP{Problème de géométrie -- production 6}

\boussole

\correctioneclair{%
\(\widehat{ACB}=40^\circ\),
\(\widehat{BCD}=60^\circ\),
puis
\(40^\circ+60^\circ+80^\circ=180^\circ\).
Donc \(A,C,E\) sont alignés et
\(AE=6{,}4\) cm.
}

\missionouverte{%
Observe et annote cette production.
Écris ensuite tes remarques.
}

\begin{center}

\etiquette{Production 6}

\vspace{0.3em}

\copieencadree[
  width=0.92\linewidth,
  height=0.60\textheight,
  keepaspectratio
]{production 6}

\end{center}

\remarques{3cm}


% =========================================================
% PAGE 5
% ALLUMETTES : RÉPONSES 2 ET 4
% =========================================================

\newpage

\titreAP{Problème des allumettes -- réponses 2 et 4}

\boussole

\noindent
\begin{minipage}[c]{0.44\linewidth}

\centering

\copieencadree[
  width=0.96\linewidth,
  height=0.17\textheight,
  keepaspectratio
]{allumettes maisons}

\end{minipage}
\hfill
\begin{minipage}[c]{0.53\linewidth}

\small

On construit des maisons avec des allumettes comme sur la figure.

Il fallait trouver le nombre d'allumettes aux étapes
4, 10 et 2007, puis donner une méthode valable pour
une \textbf{étape quelconque}.

\end{minipage}

\correctioneclair{%
À l'étape \(n\),
on peut écrire
\[
5+4(n-1)=4n+1.
\]
Donc :
étape 4 : \(17\) ;
étape 10 : \(41\) ;
étape 2007 : \(8029\).
}

\missionouverte{%
Observe et annote les deux réponses.
Écris ce que tu remarques.
}

% ---------- Réponse 2 ----------

\etiquette{Réponse 2}

\vspace{0.25em}

\begin{center}

\copieencadree[
  width=0.92\linewidth,
  height=0.13\textheight,
  keepaspectratio
]{reponse 2}

\end{center}

\vspace{0.5em}

% ---------- Réponse 4 ----------

\etiquette{Réponse 4}

\vspace{0.25em}

\begin{center}

\copieencadree[
  width=0.92\linewidth,
  height=0.14\textheight,
  keepaspectratio
]{reponse 4}

\end{center}

\remarques{3.4cm}


% =========================================================
% PAGE 6
% ALLUMETTES : RÉPONSES 6, 7 ET 8
% =========================================================

\newpage

\titreAP{Problème des allumettes -- autres réponses}

\boussole

\correctioneclair{%
Une expression correcte est
\(\boxed{4n+1}\),
où \(n\) désigne le numéro de l'étape.
}

\missionouverte{%
Observe et annote les trois réponses.
Écris ce que tu remarques.
}

% ---------- Réponses 6 et 7 ----------

\noindent
\begin{minipage}[t]{0.47\linewidth}

\etiquette{Réponse 6}

\vspace{0.3em}

\centering

\copieencadree[
  width=0.94\linewidth,
  height=0.20\textheight,
  keepaspectratio
]{reponse 6}

\end{minipage}
\hfill
\begin{minipage}[t]{0.50\linewidth}

\etiquette{Réponse 7}

\vspace{0.3em}

\centering

\copieencadree[
  width=0.94\linewidth,
  height=0.20\textheight,
  keepaspectratio
]{reponse 7}

\end{minipage}

\vspace{0.9em}

% ---------- Réponse 8 ----------

\begin{center}

\begin{minipage}{0.80\linewidth}

\etiquette{Réponse 8}

\vspace{0.3em}

\centering

\copieencadree[
  width=0.94\linewidth,
  height=0.25\textheight,
  keepaspectratio
]{reponse 8}

\end{minipage}

\end{center}

\remarques{3.7cm}

\end{document}